% Abhakseite „Das kann ich“
\eng
\einheitenkopf*{\hypertarget{abhaken}{}Das kann ich}
\setlength{\parskip}{2pt}
\abhakgruppe{Das kennst du schon}
\abhak{1}{Ich kann mit Kommazahlen multiplizieren und dividieren und das Ergebnis runden.}
\abhak{2}{Ich kann Umfang und Flächeninhalt unterscheiden und die richtige Einheit angeben.}
\abhak{3}{Ich kann eine Formel nach einer anderen Größe umstellen.}
\abhak{4}{Ich kann Zahlen quadrieren und die Quadratwurzel ziehen.}
\abhak{5}{Ich finde den Fehler beim Quadrieren.}
\abhak{6}{Ich kann mit rechtem, gestrecktem und vollem Winkel rechnen.}
\abhak{7}{Ich kann einen Anteil als Bruch und in Prozent angeben.}
\abhakgruppe{Einheit 1 von 3 · Kreisumfang}
\abhak{1}{Ich erkenne, ob Radius oder Durchmesser gegeben ist.}
\abhak{2}{Ich erkenne, ob nach dem Rand oder nach der Fläche gefragt ist.}
\abhak{3}{Ich kann einen Kreis zeichnen und Mittelpunkt, Radius und Durchmesser einzeichnen (kennst du seit Klasse 5).}
\abhak{4}{Ich kann aus dem Radius den Durchmesser berechnen und umgekehrt.}
\abhak{5}{Ich kann herausfinden, wie oft der Durchmesser in den Umfang passt.}
\abhak{6}{Ich kann den Umfang eines Kreises berechnen.}
\abhak{7}{Ich kann aus dem Umfang den Durchmesser und den Radius berechnen.}
\abhak{8}{Ich finde den Fehler beim Kreisumfang.}
\abhak{9}{Ich kann begründen, warum $u : d$ bei jedem Kreis gleich ist.}
\abhak{10}{Ich kann ausrechnen, wie weit ein Rad rollt.}
\abhakgruppe{Einheit 2 von 3 · Kreisfläche}
\abhak{11}{Ich erkenne, welche Formel ich brauche.}
\abhak{12}{Ich kann den Flächeninhalt eines Kreises berechnen.}
\abhak{13}{Ich kann den Flächeninhalt von Halbkreis und Viertelkreis berechnen.}
\abhak{14}{Ich kann aus einer Größe eines Kreises die anderen berechnen.}
\abhak{15}{Ich kann aus dem Flächeninhalt den Radius berechnen.}
\abhak{16}{Ich finde den Fehler bei der Kreisfläche.}
\abhak{17}{Ich kann unterscheiden, ob der Umfang oder der Flächeninhalt gesucht ist.}
\abhak{18}{Ich kann begründen, warum die Kreisfläche $\pi \cdot r^2$ ist.}
\abhak{19}{Ich kann mit Kreisflächen Preise vergleichen.}
\abhakgruppe{Einheit 3 von 3 · Kreisteile}
\abhak{20}{Ich erkenne, welcher Teil des Kreises grau ist.}
\abhak{21}{Ich kann den Anteil eines Kreisausschnitts aus dem Mittelpunktswinkel berechnen.}
\abhak{22}{Ich kann aus dem Anteil den Mittelpunktswinkel berechnen.}
\abhak{23}{Ich kann Bogenlänge, Flächeninhalt und Umfang eines Kreisausschnitts berechnen.}
\abhak{24}{Ich kann den Flächeninhalt eines Kreisrings berechnen.}
\abhak{25}{Ich finde den Fehler beim Anteil eines Kreisausschnitts.}
\abhak{26}{Ich kann begründen, welcher Teil eines Kreises zu einem Mittelpunktswinkel gehört.}
\abhak{27}{Ich kann mit einem Kreisausschnitt eine Fläche planen.}
